% !TEX root = IE3_expG2_prelab.tex
\documentclass[lualatex,a4paper,10pt]{jlreq}
% Shared LuaLaTeX preamble.
\usepackage{luatexja-fontspec}
\setmainfont{TeX Gyre Termes}
\setsansfont{TeX Gyre Heros}
\setmonofont{Latin Modern Mono}
\setmainjfont[BoldFont=HaranoAjiGothic-Medium]{HaranoAjiMincho-Regular}
\setsansjfont[BoldFont=HaranoAjiGothic-Bold]{HaranoAjiGothic-Medium}
\usepackage[top=22mm,bottom=22mm,left=23mm,right=23mm]{geometry}
\usepackage{amsmath,amssymb,booktabs,array,longtable,tabularx}
\usepackage{graphicx,xcolor,tikz,circuitikz}
\usetikzlibrary{arrows.meta,positioning,calc}
\usepackage{enumitem,fancyhdr,listings,needspace}
\usepackage[unicode,colorlinks=true,linkcolor=labblue,urlcolor=labteal]{hyperref}
\usepackage{bookmark}
\definecolor{labblue}{RGB}{30,70,112}
\definecolor{labteal}{RGB}{0,100,105}
\definecolor{labgray}{gray}{0.95}
\ModifyHeading{section}{font={\large\sffamily\bfseries\color{labblue}}}
\ModifyHeading{subsection}{font={\normalsize\sffamily\bfseries\color{labteal}}}
\setlength{\parindent}{1\zw}
\setlength{\parskip}{3pt}
\setlength{\emergencystretch}{2em}
\renewcommand{\arraystretch}{1.35}
\setlist{itemsep=3pt,topsep=4pt,leftmargin=2\zw}
\newcolumntype{Y}{>{\raggedright\arraybackslash}X}
\newcolumntype{C}[1]{>{\centering\arraybackslash}p{#1}}
\pagestyle{fancy}
\fancyhf{}
\fancyhead[L]{\small\color{labblue} IE3 後期・電子工学実験}
\fancyhead[R]{\small テーマ\LabCode}
\fancyfoot[C]{\thepage}
\setlength{\headheight}{16pt}
\DeclareOldFontCommand{\rm}{\normalfont\rmfamily}{\mathrm}
\lstset{basicstyle=\small\ttfamily,columns=fullflexible,keepspaces=true,
 breaklines=true,frame=single,backgroundcolor=\color{labgray},
 numbers=left,numberstyle=\tiny,showstringspaces=false,aboveskip=8pt,belowskip=8pt}
\newcommand{\blank}{\rule{22mm}{0.3pt}}
\newcommand{\answerline}{\par\noindent\rule{\linewidth}{0.25pt}\par}
\newcommand{\writing}[1]{\par\Needspace{\dimexpr#1+5mm\relax}\noindent%
 \fcolorbox{labblue!35}{white}{\parbox[c][#1][t]{\dimexpr\linewidth-2\fboxsep-2\fboxrule\relax}{\vspace{2mm}\noindent\strut\hfill\par}}\par}
\newcommand{\hint}[1]{\par\smallskip\noindent{\sffamily\bfseries\color{labteal}記述の視点：}#1\par}
\newcommand{\note}[1]{\par\smallskip\noindent\fcolorbox{labblue!45}{labblue!5}{%
 \parbox{\dimexpr\linewidth-2\fboxsep-2\fboxrule\relax}{#1}}\par\smallskip}
\newcommand{\eqerror}{\varepsilon=\frac{x_{\mathrm{meas}}-x_{\mathrm{th}}}{x_{\mathrm{th}}}\times100\ \%}
\newcommand{\LabSource}{徳山工業高等専門学校 情報電子工学科，
 『2026年度 電子工学実験（3年後期）テキスト』，テーマ\LabCode，
 \expandafter\path\expandafter{\SourcePdf}。}
\newcommand{\reporttitle}{
 \noindent{\small\sffamily\color{labteal}実験報告書・記入ガイド}\par
 \noindent{\LARGE\sffamily\bfseries \LabCode\quad\LabTitle}\par\smallskip
 \noindent 氏名：\blank\quad 班：\blank\quad 実験日：\blank\par
 \note{目的・原理・手順を確認し、実測値と自分の考察を記入する。
 考察のガイドは、検討する事象・使う測定結果・記述の方針を示す。}
}
\newcommand{\prelabtitle}{
 \noindent{\small\sffamily\color{labteal}事前学習・前提知識プリント}\par
 \noindent{\LARGE\sffamily\bfseries \LabCode\quad\LabTitle}\par\smallskip
 \noindent 氏名：\blank\quad 班：\blank\quad 予習日：\blank\par
 \note{用語、回路の動き、計算、測定表への記入の順に読む。
 例題の値は説明用であり、実験結果ではない。}
}
\newcommand{\tocrow}[3]{\hyperref[#1]{#2}& #3 &\pageref*{#1}\\[2mm]}
\newcommand{\documentmap}[1]{
 \section*{目次と概要}
 \noindent 青色の項目をクリックすると該当箇所へ移動する。\par
 {\small\begin{longtable}{@{}p{0.29\linewidth}p{0.61\linewidth}r@{}}
 \toprule {\color{labblue}\bfseries 項目} & {\color{labblue}\bfseries 読むと分かること} & 頁\\\midrule
 \endhead #1\bottomrule
 \end{longtable}}
 \clearpage
}
\newenvironment{terms}{\par\smallskip\begin{longtable}{@{}p{0.23\linewidth}p{0.73\linewidth}@{}}
 \toprule {\color{labteal}\bfseries 用語}&{\color{labteal}\bfseries 意味・回路での役割}\\\midrule\endhead}
 {\bottomrule\end{longtable}}
\newcommand{\term}[2]{\textbf{#1}&#2\\[2mm]}
\newcommand{\guidepoint}[4]{
 \Needspace{32mm}\subsection{#1}
 \noindent{\bfseries\color{labteal}考える事象：}#2\par
 \noindent{\bfseries\color{labteal}参照する結果：}#3\par
 \noindent{\bfseries\color{labteal}記述の方針：}#4\par
}
\newcommand{\reportclosing}[1]{
 \clearpage\section{結論のまとめ方}\label{sec:closing}
 \hint{#1}
 \ConclusionGuide
 \begin{enumerate}
 \item 目的に対応する最も重要な実測結果を、測定条件と数値を添えて述べる。
 \item 原理と一致した点・一致しなかった点を示す。原因は確認済みの事実と推測を分ける。
 \item 実施範囲で言えることと、未確認の条件を一文ずつまとめる。
 \end{enumerate}
 \note{書き出しの型：「○○の条件で△△を測定したところ、□□となった。
 これは◇◇と整合する。ただし××は確認しておらず、一般化には追加測定が必要である。」
 空欄は自分の実測値と根拠で埋める。}
 \noindent 結論（実験後に記入）：\writing{30mm}
 \noindent 改善案・次に確認したい条件：\writing{16mm}
 \section*{参考文献}\label{sec:references}
 \begin{enumerate}
 \item \LabSource
 \item 使用したデータシート・書籍・ウェブ資料（著者、題名、該当箇所、URL、参照日）を追加する。
 \end{enumerate}
}
\newcommand{\prelabclosing}{
 \section*{準備・出典}\label{sec:prelabsource}
 \begin{itemize}[nosep]
 \item 資料、筆記具、記録用紙、電卓と、実験条件・機器設定の記録欄。
 \item 確認問題の解答と、分からない用語・回路点への具体的な質問。
 \end{itemize}
 {\small\noindent\LabSource\par}
}
\newcommand{\simpleflow}[1]{
 \begin{center}
 \begin{tikzpicture}[node distance=5mm and 5mm,
 every node/.style={font=\small},box/.style={draw=labblue,fill=labblue!4,rounded corners=1mm,
 align=center,minimum height=10mm,text width=30mm}]
 #1
 \end{tikzpicture}
 \end{center}
}

\newcommand{\LabCode}{G2}
\newcommand{\LabTitle}{マイコンによる7セグメントLED表示制御}
\newcommand{\SourcePdf}{IE3_expG2_A4.pdf}
\begin{document}
\prelabtitle
\documentmap{
\tocrow{sec:auto1}{1 用語を一つずつ理解する}{言葉の意味、単位、回路や測定での役割を確認する。}
\tocrow{sec:auto2}{2 数字が表示されるまでをたどる}{部品と信号の経路を追い、状態が変化する理由を理解する。}
\tocrow{sec:auto3}{3 数字とセグメント符号}{数値を点灯ビット列へ変換する理由を読む。}
\tocrow{sec:auto4}{4 桁選択の負論理}{選択コードと表示する桁を対応させる。}
\tocrow{sec:auto5}{5 時間の計算}{命令ステートを合計し更新間隔を求める。}
\tocrow{sec:auto6}{6 カウンタの設計}{通常・境界・終了時の動作を決める。}
\tocrow{sec:auto7}{7 測定からレポートへ：表示と時間を検証する}{実測値から計算・作図・記述へ進む順序を例題で練習する。}
\tocrow{sec:auto8}{8 確認問題}{自分で計算・説明して、実験に必要な理解を確かめる。}
\tocrow{sec:auto9}{解答の目安}{数値と理由を照合し、単位や論理の取り違えを見直す。}
\tocrow{sec:auto10}{設計・予習欄}{この項目の仕組みと実験で確認すべき点を整理する。}
\tocrow{sec:prelabsource}{準備・出典}{持参物と資料を確認し、具体的な質問を準備する。}
}

\section{用語を一つずつ理解する}\label{sec:auto1}
\begin{terms}
\term{7セグメント}{a〜gの七つの発光部で数字を作る表示器。小数点hは別の一ビット。数字そのものと発光部の組合せは異なる情報である。}
\term{正論理・負論理}{正論理では1が有効、負論理では0が有効。資料のセグメントは正論理、桁選択は負論理なので同じ1でも役割が違う。}
\term{表示符号・表引き}{数字に対応する点灯パターンのビット列。数字を配列の番号として、対応コードを取り出す操作が表引きである。}
\term{時分割表示}{桁と符号を高速に切り替え、複数桁を順番に点灯する方式。十分速い切替えで連続した表示に見えるが、遅ければちらつく。}
\term{ステート}{命令の実行に必要なクロック単位。命令数が同じでも実行時間は同じとは限らず、各命令のステートを足す。}
\term{分岐とオーバーヘッド}{条件による実行経路の切替えと、計時そのもの以外の表示・呼出し・判定等の時間。待ちだけを数えると更新が遅れる。}
\term{境界・桁あふれ}{9の次を0に戻すなど、数の範囲を超える時の処理。通常の加算とは別に仕様を決めて試験する。}
\end{terms}
\section{数字が表示されるまでをたどる}\label{sec:auto2}
\subsection{数値を符号へ変換する}
カウンタの数値2は02hだが、数字2にはa,b,d,e,gを点灯する5Bhが必要である。
数値を一つ増やしても、符号を一つ増やしただけでは次の数字にはならない。
数値はカウンタとして保持し、毎回その数字のコードを表から取り出す。
\subsection{桁とセグメントを別々に出力する}
0Chが桁、07hが発光部の符号を担当する。
全桁を非選択にし、符号を更新してから目的の桁を選ぶと、切替え中に前の数字が別の桁へ一瞬出る誤点灯を抑えやすい。
実験の左端8Eh、右端87hは配線対応に基づくので、ビットの数値順で左右を推測しない。
\subsection{表示更新と時間更新を分ける}
複数桁の点灯を巡回する短い周期と、秒の数値を一つ増やす長い周期は役割が異なる。
一秒間まったく表示を更新しない構成では、多桁の時分割表示が維持できない場合がある。発展課題では表示の巡回と計時をどう両立させるか決める。

\section{数字とセグメント符号}\label{sec:auto3}
数字2は整数02hであるが、LEDに数字2を表示する符号は5Bhである。配列から表引きして変換する。
ビット0〜6がa〜g、ビット7が小数点hに対応する。
\[
 \text{表示コード}=\text{数字のコード}\ \mathrm{OR}\ 80h
\]
とすれば小数点を追加できる。
資料の7は27hで、fを含む字形。07hで表示する形もあるが、実験では採用した字形を明記する。
\section{桁選択の負論理}\label{sec:auto4}
4桁すべて非選択の8Fhから、選びたい桁のビットを0にする。
\par\noindent\begin{tabularx}{\linewidth}{YYYYY}
\toprule &左端&左から2&左から3&右端\\\midrule
コード&8Eh&8Dh&8Bh&87h\\
下位4ビット&1110&1101&1011&0111\\\bottomrule
\end{tabularx}
複数桁を異なる数字で表示するには、桁と符号を切り替える時分割表示が必要である。符号を変える前に全桁非選択にすると、切替え時の誤点灯を抑えやすい。
\section{時間の計算}\label{sec:auto5}
\[
 t=\frac{N_{\rm states}}{2\,457\,600}.
\]
待ちサブルーチンだけでなく、呼出し・復帰、表示、加算、比較、条件分岐まで含める。
同じループでも最後の分岐はステート数が異なる可能性がある。
\clearpage
\section{カウンタの設計}\label{sec:auto6}
保持する数値を0〜9として更新し、10になったら0に戻す。セグメントのビット列を加算して数字を作ろうとしない。
発展課題では、初期値、終了条件、更新間隔、桁あふれ、小数点の役割を先に決める。
\simpleflow{
\node[box] (a) {数値を保持};
\node[box,right=of a] (b) {符号を表引き};
\node[box,right=of b] (c) {桁選択・出力};
\node[box,right=of c] (d) {待ち・数値更新};
\draw[-{Stealth}] (a)--(b);\draw[-{Stealth}] (b)--(c);\draw[-{Stealth}] (c)--(d);
}
\section{測定からレポートへ：表示と時間を検証する}\label{sec:auto7}
\subsection{表示を一桁ずつ調べる}
数字0〜9の符号を出して、実際に点いたセグメントを記録する。7の27hではfも点く字形なので、07hとの違いを故障と決めつけない。
次に同じ符号のまま桁コードだけを変えて左右を確認する。この二段階で、符号の誤りと桁の誤りを分けられる。
\subsection{時間の足し算}
仮に待ちが$2\,457\,600$ステート、表示と判定が追加で2400ステートなら、
\[
 t=\frac{2\,460\,000}{2\,457\,600}\simeq1.00098\,\mathrm{s}.
\]
待ちだけが1秒でも、更新は約0.98 ms余分にかかる。これは説明用のステート数で、実際には命令表を参照して全経路を数える。
ループの最終分岐が通常と違う場合、回数を一つ分けて計算する。
\subsection{仕様から試験を作る}
10秒カウンタなら0→1と9→0、発展課題なら終了時、初期化直後、桁あふれを確認する。
実測時間をどの二イベントの間で測ったか示す。数回の更新をまとめて測り平均する方法では、何回分かも記録する。

\clearpage
\section{確認問題}\label{sec:auto8}
\begin{enumerate}
\item 数字0の3Fhに小数点を付けると何hになるか。\answerline
\item 右端のみを選ぶコードと、全桁を消すコードを答える。\answerline
\item $1\,228\,800$ステートの処理時間を求める。\answerline
\item 待ち処理を1秒に合わせた後に表示処理を追加すると、更新間隔はどうなるか。\answerline
\end{enumerate}
\section*{解答の目安}\label{sec:auto9}
(1) BFh。(2) 87h、8Fh。(3) 0.5秒。(4) 追加処理分だけ1秒より長くなる。
\section*{設計・予習欄}\label{sec:auto10}
発展課題の仕様と、通常状態・境界状態・終了状態の試験を三つ以上考える。
\writing{26mm}
\prelabclosing
\end{document}
